% 【专题研讨】所属：第8章（由 chapters/revised/08-digital-ai.tex \input 引入）
\minitopic{专题研讨：生成式系统、来源证明与认识责任}数字环境没有取消传统认识问题，而是重新排列了证据链。搜索排序决定哪些材料先被看见，生成系统把多种训练痕迹压缩成没有明确出处的文本，合成媒体让“看起来像记录”不再自动支持真实发生。用户面对的不是简单的“机器会不会说谎”，而是四个可以分别失败的环节：内容是否为真、输出由何过程产生、来源能否追溯、行动责任由谁承担。

\minitopic{工具、说者还是混合认识源}把生成系统当作计算器，会强调它执行输入到输出的功能，认识责任主要落在用户验证；把它当作证言者，会追问系统是否具有断言意图、能力和责任；把它当作信息环境的一部分，则关注训练、部署、检索和界面共同构成的流程。三种模型对不同任务各有解释力，不能仅凭系统使用第一人称语言就裁决。 证言模型面临主体性问题。传统证言不仅产生句子，还涉及说者把某内容呈交为真、可被质询并承担信誉后果。生成系统可以模仿这些语用标志，却未必拥有稳定信念、承诺或自主目的。工具模型则可能过度简化：用户常依赖系统选择资料和组织答案，输出对信念的影响类似社会来源，并非中性尺子。 更实用的混合分析把责任分层。开发与部署机构负责已知性能边界、数据治理和界面提示；集成者负责检索源、权限与日志；专业用户负责把输出放进合适任务并核验；最终决策者负责行动。是否称机器“知道”可以继续争论，但不能让词语争论替代风险分配。

\minitopic{流畅、置信与可靠性}语言流畅主要表明输出符合语言模式，不表明其中事实得到独立证据支持。细节丰富、语气坚定和格式完整会提高人的主观可信度，正因如此，流畅错误特别危险。核验时要把答案拆成原子主张：事实、引用、推断、建议与价值判断分别检查，不以一处正确为整段背书。 模型内部下一个符号的概率，不等于命题真实概率；系统在某基准上的平均准确率，也不等于当前个案置信度。界面若显示“90\% confidence”，必须说明对象、估计方法、校准集与适用范围。没有定义的置信数比不显示更误导，因为它把形式精确伪装成认识精确。 可靠性必须任务与环境相对。系统在常见英语事实问答上表现良好，不能推出它在地方史、最新法规或少数语言医学上同样可靠。版本更新、提示结构、检索库和工具调用都可能改变过程类型。用户的“它以前答得很准”只有在任务相似且环境稳定时才提供有限支持。

\minitopic{幻觉、遗漏与不可见错误}“幻觉”通常指系统生成貌似合理但无事实根据的内容，但这个宽泛标签掩盖不同故障：虚构实体、把真来源配给错误主张、错误合并人物、过时信息、无根据推断和遗漏关键反例。修复方法不同。存在性错误可查权威目录，引用错配需阅读原文，过时性需核对日期，遗漏则需要独立检索而非只核对现有清单。 最危险错误未必最离奇。虚构论文容易在数据库中被发现；真实论文、真实作者和真实期刊被组合成错误论点，更容易通过表面核验。用户不能止于“链接能打开”，还要检查页面是否支持紧邻主张、研究对象与结论范围是否匹配、版本是否已撤回。 系统拒答与保守表达也可能产生错误。若模型在少数群体、争议历史或低资源语言上更常说“不确定”，信息缺口会分布不均。质量审计应同时测量错误断言、错误拒答与重要遗漏。只追求降低虚构率可能让系统变得无用，合理目标是按风险校准回答、来源和拒答。

\minitopic{来源可追溯不等于内容为真}内容来源标准可以记录媒体的创建、编辑与签名历史，例如为数字资产附加可验证的来源声明。\parencite{c2pa2025} 这类机制回答“文件经历了什么”，不能单独回答“文件所描绘事件是否真实”。一张有完整相机签名的照片仍可能摆拍，一篇来源清楚的文章仍可能推理错误；反过来，缺少签名的旧照片也不因此必假。 来源证明的认识价值在于缩小故障空间。签名链完整可排除某些未经记录的修改，原始设备和时间信息能支持进一步调查；链条断裂则提示不确定，但不能自动判伪。验证界面应展示具体断裂位置，而不是只给绿色或红色徽章。 制度还要考虑密钥泄露、错误元数据、平台剥离和采用不均。记者、档案馆与普通用户的工作流不同，强制来源标记可能无意中降低未采用技术者的可信度。技术标准应与内容核查、来源访谈和异质记录结合，不把真实性外包给单一协议。

\minitopic{检索增强与引用核验}带检索的生成系统可以把答案连接到可查看文档，减少完全无来源输出，却不保证引用充分。系统可能检到相关但不蕴涵主张的页面，只引用摘要忽略限制，或把多个来源拼接成任何一篇都未支持的结论。引用质量应按“存在、相关、蕴涵、范围、版本”五项检查。 检索库本身也有边界。企业知识库可能缺少最新更正，网页搜索受排序与可访问性影响，学术数据库覆盖语言和学科不均。系统没有检到反例，不等于没有反例。高风险任务需要说明检索范围、日期、关键词和排除标准，并用独立路径复查关键主张。 好的界面应把句子与来源片段对齐，区分原文事实和系统推断，允许用户查看上下文。用户则要抵抗“有脚注即可靠”的形式偏差。引用是一条待验证的证据边，不是装饰性印章。

\minitopic{自动化偏差与有意义的人类监督}“人工在环”只有在人工拥有时间、能力、信息和否决权时才有认识意义。操作员每小时确认数千条警报，或只能在接受按钮旁看到系统结论，监督只是责任转移。有效安排会优先呈现不确定和分歧案例，保存人工初判，要求覆写理由并抽样复核机器与人一致的案例。 也要防止算法厌恶。一项显眼机器错误可能让人放弃整体较可靠工具，而对人类同类错误视而不见。比较应使用相同任务、相同损失和分群数据，问人机错误是否互补。若两者总在同类困难案例失败，简单投票不会提高表现。 人机协作的目标不是最大化接受率，而是改进最终决定与错误发现。可以把系统分配给高召回筛查、让人类处理语境与价值；也可让人先独立判断，再用系统作为第二意见。哪种结构更好，需要在真实工作负荷和环境变化下验证。

\minitopic{教育与研究中的可审计使用}在学习中，生成系统可用于提出反例、改变案例变量、检查论证遗漏和提供表达反馈。若直接替代问题重构与资料阅读，它会制造“看过答案”的熟悉感而非理解。一个稳健流程是：先独立写出立场和不确定点，再请求系统给最强反对，逐项核验事实，最后说明哪些理由改变了原判断。 研究使用应保存工具名称、版本或访问日期、关键输入、检索范围、采用与拒绝的输出，以及人工核验来源。记录不是把系统列为作者，而是让工作链可复现。涉及个人资料、未公开研究或受版权限制文本时，还要先确定是否有权上传。 最终文章只能引用可核对来源，不能把“模型说”当作学术证据。系统可以帮助发现文献，却不代替阅读；可以建议论证，却不承担作者对断言的责任。披露规则随机构变化，但无论是否要求公开披露，内部可追溯性都应保留。

\minitopic{公平、分布变化与风险治理}平均性能会掩盖分布性失败。面部分析研究曾展示商业系统在交叉群体上的错误率差异，说明审计不能只给总体数字。\parencite{buolamwini2018} 公平也不止一种指标：错误率平衡、校准、机会与个体相似性可能在不同基础率下冲突。选择指标包含价值判断，必须结合系统用途与受害路径。 部署后环境会变化。新疾病、新语言用法、政策改变或用户策略都可能使训练评估失效。风险治理应持续监测而非一次认证：记录漂移信号、定义停用阈值、允许受影响者申诉、定期重新评估并明确谁能暂停系统。NIST 的生成式 AI 风险管理资料强调将风险识别、测量与管理嵌入完整生命周期。\parencite{nist2024} 公平修复也可能产生新误差。为一个群体调阈值会改变假阳性与假阴性分配，增加数据收集又带来隐私风险。报告应列出受影响者、可逆性和剩余风险，不以单个“负责任 AI”标签遮蔽冲突。

\minitopic{综合案例：自动生成的法律摘要}法院内部试用系统，把长篇案卷生成供法官预览的摘要。系统给出的当事人姓名和时间线大体正确，却把证人“未能确认看见被告”压缩成“证人确认被告在场”。摘要附有案卷链接，但没有定位页码。法官助理因工作量大，只抽查开头。 分析要区分事实错误、否定范围丢失、来源粒度不足和监督能力不足。总体摘要准确率无法表示这类关键语义反转的风险。系统若只供导航，错误仍可能通过注意锚定影响后续阅读；标注“不得依赖”不能消除界面因果作用。 改进包括逐句定位、对否定与不确定表达设置保守规则、先由助理独立列争点、对影响权利的命题强制原文复核、抽样审计一致与不一致案例，并保留更正记录。若无法提供足够时间和原文可达性，就不应把“人工在环”当作安全依据。最终责任需要机构、产品团队、助理和法官按可控环节分配，而非全部压给最后点击者。

\begin{tcolorbox}[breakable,colback=white,colframe=blue!30!black,title={专题研读任务},fonttitle=\bfseries]
选取一个生成式系统完成同一事实任务十次。建立错误类型表，分别记录虚构、错引、过时、遗漏、范围夸大和拒答；再设计一个来源可追溯但内容仍错误的案例，说明来源证明究竟排除了什么、没有排除什么。
\end{tcolorbox}

\index{生成式人工智能}
\index{幻觉}
\index{来源可追溯性}
\index{检索增强生成}
\index{人工在环}
\index{自动化偏差}
\index{分布变化}
